% =====================================================================
% sections/background.tex — long version only (~2-3 pages)
% =====================================================================

\section{Institutional background}
\label{sec:background}

\subsection{The 1989 program launch}

The Village Midwife Program (\emph{Bidan di Desa}) was announced in
1989 in response to a series of national maternal-mortality reviews
that had identified the absence of skilled birth attendants in rural
villages as the single largest contributor to Indonesia's high
maternal mortality ratio \autocite{geefhuysen1999,gani1996}. The
program had three structural features that shape the empirical design
in this paper.

First, midwives were trained in formal three-year midwifery programs
attached to provincial nursing academies and were certified as full
allied health professionals---not as paraprofessional health workers
of the kind deployed by many similar programs across South and
Southeast Asia. Second, each midwife was deployed to live in the
village she served, with a guaranteed government salary for at least
the first three years. Third, after the contract period the midwife
was permitted to remain in private practice in the same community,
with the expectation that she would have built a sufficient client
base over the contract period to sustain the practice on
out-of-pocket payments \autocite{gani1996}. By 1996, approximately
\num{54000} village midwives had been deployed
\autocite{frankenberg2005}.

The pace and geography of rollout reflected three pragmatic
constraints. The supply of newly trained midwives was bounded by the
capacity of provincial nursing academies, which expanded throughout
the 1990s. The deployment process was organized by provincial health
departments, which prioritized communities with the weakest existing
maternal and child health infrastructure---fewer existing facilities,
greater distance to the nearest puskesmas, and higher baseline
mortality. The 1997 Asian Financial Crisis disrupted the rollout: the
guaranteed-salary contract was difficult to renew during fiscal
contraction, and many midwives whose initial three-year contracts
expired in the 1997--2000 window left their assigned villages for
urban private practice or family responsibilities
\autocite{triyana2016,jointcommittee2013}. The contraction in the
midwife stock during the AFC window introduces a second source of
variation in village-level access to the program; the relevant data
for measuring it are unfortunately too sparse to support a precise
analysis (the SAR's exit-year fields are not consistently filled), so
the present paper focuses on the entry shock.

\subsection{Services provided}

The scope of services provided by village midwives extended well
beyond skilled birth attendance, even though delivery care was the
program's stated rationale. \textcite{makowiecka2008} document the
range of services in two representative districts: antenatal care
visits with anthropometric measurement and counseling on diet,
postnatal home visits in the first month after delivery, family
planning consultations and contraceptive distribution, immunization
referral, and primary preventive care for non-pregnant women and
children. Village midwives also distributed iron supplements to
pregnant women, vitamin~A to postnatal mothers and infants, and
oral rehydration salts to children with diarrheal illness. The
breadth of the service bundle is the primary reason the program is
naturally interpreted as an early-life intervention rather than a
narrowly defined obstetric service.

\subsection{Non-random rollout and its consequences}

The deliberate prioritization of communities with weaker baseline
infrastructure means that the year a community received a village
midwife, $\text{start}_c$, is mechanically negatively correlated with
village-level maternal and child health outcomes that the program was
designed to improve. Communities that received the midwife earliest
were the same communities where a child born at the same calendar
time would have been at greatest health risk in the absence of the
program. A naïve cross-sectional comparison of cohorts treated early
to those treated late would therefore conflate the program's effects
with the consequences of the underlying community vulnerability that
drove placement.

The empirical strategy in Section~\ref{sec:empirical} responds to this
problem in three ways. The TWFE specification absorbs all
time-invariant village characteristics into community-of-birth fixed
effects. The within-mother specification additionally absorbs every
village characteristic that is time-invariant during the relevant
cohort window, observable or not, by comparing siblings born to the
same mother in the same village. And the
\textcite{callawaysantanna2021} estimator uses never-treated
communities (those that never received a midwife within the
observation window) as a control group, sidestepping the staggered-DiD
biases that arise when already-treated communities are used as
controls for later-treated ones \autocite{goodmanbacon2021}.

These corrections do not eliminate the problem. Time-varying village
shocks correlated with rollout year---for example, the simultaneous
arrival of other facility types or the construction of paved roads
linking a village to its district capital---will still confound the
estimate. The robustness suite in Section~\ref{sec:robustness}
includes controls for the cohort-specific arrival of doctors and
clinics in each village. The fundamental point, however, is one of
emphasis: the empirical strategy here is built around the
expectation that placement bias is severe rather than around the hope
that observable controls can cure it.
